% GENERATED FILE. Edit canonical lessons, then run npm run generate:books.
% canonical-source-hash: fnv1a64:98585d928c904bab
% canonical-lessons: JA-C332-kenri, JA-C332-gimu, JA-C332-byoudou, JA-C332-sabetsu, JA-C332-zen-in

\chapter{Rights and Duties}
\label{ch:ja-rights-and-duties}
\hlchaptermodality{332}

\begin{chapteropening}
\textbf{By the end of this chapter:} \emph{I can say what people have a right to and what they must do.}

Complete the last lesson of chapter 332: I can say what people have a right to and what they must do.
\end{chapteropening}

\section[\texorpdfstring{\ja{けんり} (kenri)}{けんり (kenri)}]{\ja{けんり} (kenri) — a right}

\label{lesson:JA-C332-kenri}



\begin{quote}
Before the new one: say the Japanese for politics, then the Japanese for a law.
\end{quote}

\subsection*{You'll want to know: \ja{けんり}}
\textbf{\ja{けんり}} — \emph{kenri} — \textquotedblleft{}a right\textquotedblright{}.

\begin{quote}
\textbf{\ja{こどもにも} \ja{けんりが} \ja{あると} \ja{おもいます。}} — \emph{kodomo ni mo kenri ga aru to omoimasu} — \textquotedblleft{}I think children have rights too.\textquotedblright{}
\end{quote}

\begin{cousinweb}
\textbf{\ja{けんり}} (\emph{kenri}) is Sino-Japanese (kango): two Chinese characters read with their on'yomi, readings borrowed from Chinese. They are \emph{ken}, \textquotedblleft{}rights\textquotedblright{} (Middle Chinese \emph{gjwen}) and \emph{ri}, \textquotedblleft{}advantage\textquotedblright{} (\emph{lijH}) (KANJIDIC2; Middle Chinese in Baxter's transcription).
\end{cousinweb}

\begin{grammarlens}[title={What you've built}]
One more word for giving reasons and opinions.
\end{grammarlens}

\subsection*{Your turn}
\begin{itemize}
\raggedright
  \item \emph{Say it:} \emph{kenri}
  \item \emph{Say it:} \emph{kenri}, once more
  \item \emph{Say it:} \emph{kodomo ni mo kenri ga aru to omoimasu} — I think children have rights too.
  \item \emph{Recall:} say the Japanese for politics, then the Japanese for a law, then say \emph{kenri} again
\end{itemize}

\begin{tcolorbox}[breakable,colback=teal!4,colframe=teal!35!black,title={Before you move on}]
What does \ja{けんり} mean? (\textquotedblleft{}A right\textquotedblright{}.) Say it once more.
\end{tcolorbox}
\section[\texorpdfstring{\ja{ぎむ} (gimu)}{ぎむ (gimu)}]{\ja{ぎむ} (gimu) — a duty; an obligation}

\label{lesson:JA-C332-gimu}



\begin{quote}
Before the new one: say the Japanese for a law, then the Japanese for a right.
\end{quote}

\subsection*{You'll want to know: \ja{ぎむ}}
\textbf{\ja{ぎむ}} — \emph{gimu} — \textquotedblleft{}a duty; an obligation\textquotedblright{}.

\begin{quote}
\textbf{\ja{しみんには} \ja{ぎむも} \ja{あると} \ja{おもいます。}} — \emph{shimin ni wa gimu mo aru to omoimasu} — \textquotedblleft{}I think citizens have duties as well.\textquotedblright{}
\end{quote}

\begin{cousinweb}
\textbf{\ja{ぎむ}} (\emph{gimu}) is Sino-Japanese (kango): two Chinese characters read with their on'yomi, readings borrowed from Chinese. They are \emph{gi}, \textquotedblleft{}righteousness\textquotedblright{} and \emph{mu}, \textquotedblleft{}duties\textquotedblright{} (Middle Chinese \emph{mjuH}) (KANJIDIC2; Middle Chinese in Baxter's transcription).
\end{cousinweb}

\begin{grammarlens}[title={What you've built}]
One more word for giving reasons and opinions.
\end{grammarlens}

\subsection*{Your turn}
\begin{itemize}
\raggedright
  \item \emph{Say it:} \emph{gimu}
  \item \emph{Say it:} \emph{gimu}, once more
  \item \emph{Say it:} \emph{shimin ni wa gimu mo aru to omoimasu} — I think citizens have duties as well.
  \item \emph{Recall:} say the Japanese for a law, then the Japanese for a right, then say \emph{gimu} again
\end{itemize}

\begin{tcolorbox}[breakable,colback=teal!4,colframe=teal!35!black,title={Before you move on}]
What does \ja{ぎむ} mean? (\textquotedblleft{}A duty; an obligation\textquotedblright{}.) Say it once more.
\end{tcolorbox}
\section[\texorpdfstring{\ja{びょうどう} (byōdō)}{びょうどう (byōdō)}]{\ja{びょうどう} (byōdō) — equal; equality}

\label{lesson:JA-C332-byoudou}



\begin{quote}
Before the new one: say the Japanese for a right, then the Japanese for a duty.
\end{quote}

\subsection*{You'll want to know: \ja{びょうどう}}
\textbf{\ja{びょうどう}} — \emph{byōdō} — \textquotedblleft{}equal; equality\textquotedblright{}.

\begin{quote}
\textbf{\ja{おとこと} \ja{おんなは} \ja{びょうどうです。}} — \emph{otoko to onna wa byōdō desu} — \textquotedblleft{}Men and women are equal.\textquotedblright{}
\end{quote}

A na-adjective: \ja{な} before a noun.

\begin{cousinweb}
\textbf{\ja{びょうどう}} (\emph{byōdō}) is Sino-Japanese (kango): two Chinese characters read with their on'yomi, readings borrowed from Chinese. They are \emph{byō}, \textquotedblleft{}even\textquotedblright{} (Middle Chinese \emph{bjaeng}) and \emph{dō}, \textquotedblleft{}equal\textquotedblright{} (\emph{tongX}) (KANJIDIC2; Middle Chinese in Baxter's transcription).
\end{cousinweb}

\begin{grammarlens}[title={What you've built}]
One more word for giving reasons and opinions.
\end{grammarlens}

\subsection*{Your turn}
\begin{itemize}
\raggedright
  \item \emph{Say it:} \emph{byōdō}
  \item \emph{Say it:} \emph{byōdō}, once more
  \item \emph{Say it:} \emph{otoko to onna wa byōdō desu} — Men and women are equal.
  \item \emph{Recall:} say the Japanese for a right, then the Japanese for a duty, then say \emph{byōdō} again
\end{itemize}

\begin{tcolorbox}[breakable,colback=teal!4,colframe=teal!35!black,title={Before you move on}]
What does \ja{びょうどう} mean? (\textquotedblleft{}Equal; equality\textquotedblright{}.) Say it once more.
\end{tcolorbox}
\section[\texorpdfstring{\ja{さべつ} (sabetsu)}{さべつ (sabetsu)}]{\ja{さべつ} (sabetsu) — discrimination}

\label{lesson:JA-C332-sabetsu}



\begin{quote}
Before the new one: say the Japanese for a duty, then the Japanese for equal.
\end{quote}

\subsection*{You'll want to know: \ja{さべつ}}
\textbf{\ja{さべつ}} — \emph{sabetsu} — \textquotedblleft{}discrimination\textquotedblright{}.

\begin{quote}
\textbf{\ja{さべつを} \ja{しない} \ja{ほうが} \ja{いいと} \ja{おもいます。}} — \emph{sabetsu o shinai hō ga ii to omoimasu} — \textquotedblleft{}I think it is better not to discriminate.\textquotedblright{}
\end{quote}

\begin{cousinweb}
\textbf{\ja{さべつ}} (\emph{sabetsu}) is Sino-Japanese (kango): two Chinese characters read with their on'yomi, readings borrowed from Chinese. They are \emph{sa}, \textquotedblleft{}distinction\textquotedblright{} (Middle Chinese \emph{tsrhae}) and \emph{betsu}, \textquotedblleft{}separate\textquotedblright{} (\emph{bjet}) (KANJIDIC2; Middle Chinese in Baxter's transcription).
\end{cousinweb}

\begin{grammarlens}[title={What you've built}]
One more word for giving reasons and opinions.
\end{grammarlens}

\subsection*{Your turn}
\begin{itemize}
\raggedright
  \item \emph{Say it:} \emph{sabetsu}
  \item \emph{Say it:} \emph{sabetsu}, once more
  \item \emph{Say it:} \emph{sabetsu o shinai hō ga ii to omoimasu} — I think it is better not to discriminate.
  \item \emph{Recall:} say the Japanese for a duty, then the Japanese for equal, then say \emph{sabetsu} again
\end{itemize}

\begin{tcolorbox}[breakable,colback=teal!4,colframe=teal!35!black,title={Before you move on}]
What does \ja{さべつ} mean? (\textquotedblleft{}Discrimination\textquotedblright{}.) Say it once more.
\end{tcolorbox}
\section[\texorpdfstring{\ja{ぜんいん} (zen'in)}{ぜんいん (zen'in)}]{\ja{ぜんいん} (zen'in) — everyone (in a group); all members}

\label{lesson:JA-C332-zen-in}



\begin{quote}
Before the new one: say the Japanese for equal, then the Japanese for discrimination.
\end{quote}

\subsection*{You'll want to know: \ja{ぜんいん}}
\textbf{\ja{ぜんいん}} — \emph{zen'in} — \textquotedblleft{}everyone (in a group); all members\textquotedblright{}.

\begin{quote}
\textbf{\ja{ぜんいんが} \ja{さんせいしたので、きまりました。}} — \emph{zen'in ga sansei shita node, kimarimashita} — \textquotedblleft{}Everyone agreed, so it was decided.\textquotedblright{}
\end{quote}

\begin{cousinweb}
\textbf{\ja{ぜんいん}} (\emph{zen'in}) is Sino-Japanese (kango): two Chinese characters read with their on'yomi, readings borrowed from Chinese. They are \emph{zen}, \textquotedblleft{}all\textquotedblright{} (Middle Chinese \emph{dzjwen}) and \emph{in}, \textquotedblleft{}member\textquotedblright{} (\emph{hjwen}) (KANJIDIC2; Middle Chinese in Baxter's transcription).
\end{cousinweb}

\begin{grammarlens}[title={What you've built}]
One more word for giving reasons and opinions.
\end{grammarlens}

\subsection*{Your turn}
\begin{itemize}
\raggedright
  \item \emph{Say it:} \emph{zen'in}
  \item \emph{Say it:} \emph{zen'in}, once more
  \item \emph{Say it:} \emph{zen'in ga sansei shita node, kimarimashita} — Everyone agreed, so it was decided.
  \item \emph{Recall:} say the Japanese for equal, then the Japanese for discrimination, then say \emph{zen'in} again
\end{itemize}

\begin{tcolorbox}[breakable,colback=teal!4,colframe=teal!35!black,title={Before you move on}]
What does \ja{ぜんいん} mean? (\textquotedblleft{}Everyone (in a group); all members\textquotedblright{}.) Say it once more.
\end{tcolorbox}
